\section{Introduction and motivation}
\label{sec:intro}



MINERvA (Main Injector Experiment for $\nu$-A) is a dedicated neutrino--nucleus scattering experiment at Fermilab, situated in the NuMI beamline immediately upstream of the MINOS near detector~\cite{MINERvA:2013zvz}. Over roughly a decade of operation, ending in 2019, MINERvA recorded high-statistics samples in both the low-energy and medium-energy (ME) NuMI beam configurations~\cite{MINERvA:2016iqn} and in neutrino-dominated (forward horn current, FHC) and antineutrino-dominated beam modes. Its physics goal is the precision measurement of neutrino--nucleus interaction cross sections on a range of nuclear targets in the few-GeV region most relevant to accelerator-based oscillation experiments: these cross sections, and the nuclear effects that shape them, are a leading systematic uncertainty for experiments such as T2K~\cite{T2K:2023osc}, NOvA~\cite{NOvA:2021nfi}, Hyper-Kamiokande~\cite{HyperK:2018design}, and DUNE~\cite{DUNE:2020jqi}, and the few-GeV regime mixes quasistatic, resonant, and deep-inelastic processes that no single model describes cleanly. This note uses the public ME FHC CC-inclusive data product; the detector itself is described in \S\ref{sec:experiment}.

The MINERvA ME FHC charged-current (CC) inclusive
double-differential cross section in muon transverse and longitudinal momentum,
$d^{2}\sigma/(d\pT\,d\pPar)$, was published by Ruterbories \emph{et al.}
\cite{MINERvA:2021owq} using D'Agostini iterative Bayesian unfolding
(IBU)~\cite{DAgostini:1994fjx} with a binned response matrix. This note reproduces
that measurement with \emph{unbinned} OmniFold~\cite{Andreassen:2019cjw,
Andreassen:2021zzk}, an unfolding technique based on iterated reweighting with
machine-learned likelihood ratios that operates on per-event observables rather
than on a fixed histogram binning. Recent advances in machine learning have enabled unbinned unfolding
techniques; see Refs.~\cite{Arratia:2021otl,Huetsch:2024quz} for overviews
and Ref.~\cite{unbinnedguide} for cross-experiment practical
guidance. These methods fall into two broad families: those that
iteratively reweight a starting simulation, of which
OmniFold~\cite{Andreassen:2019cjw,Andreassen:2021zzk,Chan:2023tbf} is the
most widely adopted, and those that generate new samples directly with a
neural network~\cite{Datta:2018mwd,Bellagente:2019uyp,Bellagente:2020piv,
Backes:2022vmn}. OmniFold was first demonstrated on collider data by the H1
collaboration~\cite{H1:2021wkz,H1:2023fzk} and has since been adopted across
a number of experiments; the application closest to this work is the
OmniFold study on simulated T2K near-detector data~\cite{Huang:2024omnifold}.
These techniques lift the binned-response limitations of widely-used
classical unfolding methods~\cite{Adye:2011gm} like IBU, the SVD
method~\cite{Hocker:1995kb}, and TUnfold~\cite{Schmitt:2012kp}. 



\paragraph{Why unbinned:}
\begin{itemize}
  \item \textbf{No binning at unfold time.} OmniFold learns event weights in the
        full observable space; the result is histogrammed only at the end, so the
        binning can be chosen (or refined) after unfolding without re-running.
  \item \textbf{Natural path to higher dimensions.} Adding observables is adding
        input features, not multiplying response-matrix bins --- this is the
        enabling property for the available-energy extension
        (\S\ref{sec:3d}) that has no practical IBU analog.
  \item \textbf{Different regularization} OmniFold's regularizer is the machine learning model capacity in addition to the iteration count in common with IBU. Comparing the two
        approaches on the same data quantifies the method-dependence of the result
        (\S\ref{sec:results}).
\end{itemize}

\paragraph{Analysis strategy.}
The published 2D cross section provides a quantitative validation of unbinned
OmniFold on MINERvA data, including its uncertainty budget. The validated
pipeline is then extended by adding $\Eavail$, $q_{3}$ and $W$ as classifier
inputs. Marginalization to each lower-dimensional result, closure, coverage and
classifier-family comparisons test the extensions. A recoil-only point-cloud
cross-check and a scalar study with relaxed truth-muon cuts probe alternative
representations and signal definitions.
